\documentclass[letterpaper, 10pt, conference]{ieeeconf}
\IEEEoverridecommandlockouts
\overrideIEEEmargins

\usepackage{graphics}
\usepackage{graphicx}
\usepackage{mathptmx}
\usepackage{times}
\usepackage{amsmath}
\usepackage{amssymb}
\usepackage{booktabs}
\usepackage{array}
\usepackage{multirow}
\usepackage{subcaption}
\usepackage[capitalize]{cleveref}

\title{\LARGE \bf Experiments and Results}

\begin{document}
\maketitle
\thispagestyle{empty}
\pagestyle{empty}


% ============================================================================
\section{Experiments}
\label{sec:experiments}

We evaluate the proposed SE(3) grasp-aware workspace optimization
on six tabletop manipulation scenarios of increasing difficulty.
The experiments address two research questions:

\begin{itemize}
  \item[\textbf{RQ1}] Does the proposed workspace optimization
    improve goal inference compared to unoptimized (random) layouts?
    We measure both the accuracy and speed of intent recognition
    under a Boltzmann path-efficiency observer matched to the
    optimizer's objective.
  \item[\textbf{RQ2}] How well does the optimization scale with
    respect to the number of candidate goals?  We evaluate on
    scenes ranging from $M{=}2$ to $M{=}8$ competing pick objects
    and report both inference performance and optimizer compute cost.
\end{itemize}

% ------------------------------------------------------------------
\subsection{Experimental Setup}
\label{sec:setup}

\paragraph{Scenarios.}
We construct six pick-and-return tasks on a shared Sawyer
tabletop workspace (\cref{tab:scenarios}).  Each task requires the
robot to identify the human operator's intended pick object from
$M$ candidates using early end-effector motion.  All objects remain
as candidates at every pick state; after each pick the object is
returned to its origin and the inference resets.  We evaluate every
permutation of pick orderings (capped at $120$ sampled orderings
for $M \ge 5$).

\begin{table}[t]
\centering
\caption{Evaluation scenarios.  $M$ is the number of pick candidates
at each step.  ``Total'' includes fixed non-pick objects.}
\label{tab:scenarios}
\small
\begin{tabular}{@{}llccc@{}}
\toprule
\textbf{Tier} & \textbf{Scene} & $M$ & \textbf{Total} & \textbf{Grasp mix} \\
\midrule
Easy   & breakfast\_easy  & 2 & 3  & top \\
Med-A  & desk             & 3 & 5  & side + top \\
Med-B  & breakfast        & 3 & 6  & side + top \\
Med-C  & kitchen\_prep    & 3 & 5  & side + top \\
Hard-A & meal\_assembly   & 5 & 7  & side + top \\
Hard-B & cluttered        & 8 & 11 & mixed \\
\bottomrule
\end{tabular}
\end{table}

\paragraph{Grasp library.}
Each object has a single pre-assigned SE(3) grasp pose from a
hand-authored library mixing top-down and side approaches.  Grasp
offsets are computed per-object from the actual half-extents so
the gripper fingertips contact the object surface.  The grasp
library is fixed across all conditions; only object positions and
yaws change.

\paragraph{Conditions.}
We compare three layout conditions:
\begin{enumerate}
  \item \textbf{Random} ($n{=}30$ layouts per scene): object
    $(x,y)$ positions sampled uniformly within table bounds,
    subject to minimum pairwise clearance ($12$\,cm), robot
    exclusion ($15$\,cm radius), and table-edge margins.  Object
    yaws sampled from $\mathcal{U}(-\pi, \pi)$.
  \item \textbf{DE-optimized} (single layout): differential
    evolution over $(x,y)$ with nested yaw gradient descent,
    maximizing the SE(3) separability slack subject to collision-free
    grasp feasibility.
  \item \textbf{ME-optimized} (best elite from a MAP-Elites
    archive): CMA-ME with the same SE(3) fitness pipeline but
    producing a diverse archive ($20 {\times} 20$ grid) of feasible
    layouts.  Behavior features: mean pairwise distance (spread)
    and centroid offset from the table center (asymmetry).
\end{enumerate}

\paragraph{Simulator.}
All evaluation runs use a headless MuJoCo simulator with the Sawyer
robot under full 3D Jacobian-pseudoinverse velocity control at
$20$\,Hz.  The simulated joystick produces $\mathbf{u}_t =
\boldsymbol{\pi}_h(\mathbf{x}_t, \mathbf{g}^*) +
\boldsymbol{\epsilon}_t$ with Gaussian noise $\sigma_v {=} 0.03$\,m/s
and peak speed $0.05$\,m/s.  The end-effector moves in full 3D
toward the target grasp pose.  The inference loop terminates when
the EE has stopped advancing (windowed stall detection, $5$\,mm net
displacement over $1$\,s), when the posterior crosses the commit
threshold, or at a $30$\,s timeout.

\paragraph{Observer.}
The SE(3) Boltzmann path-efficiency observer assigns probability to
each goal $g$ based on the efficiency of the observed trajectory:
$P(g \mid \xi) \propto \exp(-\beta \cdot C(\xi, g))$ where $C$ is
the path-efficiency cost ratio using the SE(3) distance metric with
rotation weight $\lambda_R {=} 0.04$.  We use $\beta {=} 5$ and a
default commit threshold $p_{\mathrm{thresh}} {=} 0.9$.

\paragraph{Feasibility check.}
Before running inference, each layout is checked for grasp
feasibility: per-object IK reachability and arm-vs-other-object
3D collision detection at the IK-solved joint configuration.
Infeasible grasps (no IK solution or arm collision with a
non-target object) are skipped during inference but counted as
failures in the task success metric.

\paragraph{Metrics.}
\begin{itemize}
  \item \textbf{Argmax accuracy}: fraction of feasible pick steps
    where the posterior mode at termination was the correct goal.
    Threshold-independent.
  \item \textbf{Time to inference}: mean time (s) from EE departure
    to commitment or stall, over all pick steps.
  \item \textbf{Task success rate}: fraction of complete task
    executions where \emph{every} pick had a feasible grasp.
    One infeasible grasp fails the entire task.
  \item \textbf{Threshold accuracy}: fraction of picks where the
    posterior first crossed $p_{\mathrm{thresh}}$ on the correct
    goal.
\end{itemize}

% ============================================================================
\section{Results}
\label{sec:results}

% ------------------------------------------------------------------
\subsection{RQ1: Inference Accuracy and Speed}
\label{sec:rq1_results}

\cref{tab:main_results} presents the headline comparison across all
six scenarios, combining the optimizer's own objective value
(trajectory-margin slack) with the downstream inference metrics
(argmax accuracy and time to inference).

\begin{table*}[t]
\centering
\caption{Trajectory-margin slack, argmax accuracy, and time to
inference across tiers.  Random denotes $10$ random feasible layouts
per tier with $30$ orderings each; DE and ME are single best-elite
layouts with up to $120$ orderings.  SA runtime: SE(3) Boltzmann
path-efficiency with $\beta{=}5$, joystick noise $\sigma_u{=}0.03$
m/s, control rate $5$ Hz, goal timeout $30$ s, arrival distance
$2$ cm, threshold $p_{\mathrm{thresh}}{=}0.9$. The SA loop commits
on the current argmax when the EE reaches within $2$\,cm of the true
target, matching a deployed system where the user physically arrives
at the goal. Slack is the worst-case pairwise trajectory-margin
slack computed with matched optimizer parameters ($\sigma_u{=}0.03$,
$dt{=}0.20$, $n_\text{steps}{=}75$ representing the expected pick
horizon of $15$\,s at $5$\,Hz, $\alpha{=}0.05$); positive slack
implies a nominal argmax-accuracy guarantee of
$\geq 1 - \alpha = 95\%$ under the straight-line, linear-Gaussian
approximation (\cref{sec:traj_margin}). Task success and
infeasibility are not shown in this table: random has
$5$--$49\%$ per-pick infeasibility leading to task success as low as
$10\%$, while every DE and ME entry achieves $100\%$ task success
with $0\%$ infeasibility.}
\label{tab:main_results}
\small
\setlength{\tabcolsep}{5pt}
\begin{tabular}{@{}llccccc@{}}
\toprule
\textbf{Tier} & \textbf{Layout} & $M$ &
  Slack &
  $T_{\mathrm{infer}}$ (s) &
  Argmax acc. \\
\midrule
\multirow{3}{*}{Easy}
  & Random    & 2 & ---             & $6.1 \pm 2.4$ & $92\%$ \\
  & DE        &   & $8.22$          & $\mathbf{4.1}$ & $\mathbf{100\%}$ \\
  & ME        &   & $\mathbf{8.40}$ & $5.4$          & $\mathbf{100\%}$ \\
\midrule
\multirow{3}{*}{Med-A}
  & Random    & 3 & ---  & $\mathbf{6.9 \pm 3.4}$ & $91\%$ \\
  & DE        &   & $5.03$          & $5.3$ & $\mathbf{100\%}$ \\
  & ME        &   & $\mathbf{6.04}$ & $9.4$ & $\mathbf{100\%}$ \\
\midrule
\multirow{3}{*}{Med-B}
  & Random    & 3 & ---             & $9.9 \pm 2.8$ & $\mathbf{100\%}$ \\
  & DE        &   & $4.80$          & $\mathbf{8.1}$ & $\mathbf{100\%}$ \\
  & ME        &   & $\mathbf{5.26}$ & $9.5$          & $\mathbf{100\%}$ \\
\midrule
\multirow{3}{*}{Med-C}
  & Random    & 3 & ---             & $7.7 \pm 3.9$ & $93\%$ \\
  & DE        &   & $6.00$          & $\mathbf{3.6}$ & $\mathbf{100\%}$ \\
  & ME        &   & $\mathbf{6.33}$ & $5.1$          & $\mathbf{100\%}$ \\
\midrule
\multirow{3}{*}{Hard-A}
  & Random    & 5 & ---             & $12.5 \pm 3.0$ & $88\%$ \\
  & DE        &   & $0.58$          & $18.0$         & $94\%$ \\
  & ME        &   & $\mathbf{2.28}$ & $\mathbf{10.6}$ & $\mathbf{100\%}$ \\
\midrule
\multirow{3}{*}{Hard-B}
  & Random    & 8 & ---             & $10.9 \pm 3.6$ & $25\%$ \\
  & DE        &   & $-1.89$         & $17.6$         & $22\%$ \\
  & ME        &   & $\mathbf{0.00}$ & $\mathbf{9.6}$ & $\mathbf{54\%}$ \\
\midrule
\multicolumn{2}{l}{\textit{mean across tiers}} & & \textit{DE: 3.79 / ME: 4.72} & & \textit{Random: 82\% / DE: 86\% / ME: 92\%} \\
\bottomrule
\end{tabular}
\end{table*}

\paragraph{Slack.}
Both DE and ME trajectory-margin slacks decrease with difficulty,
with ME consistently above DE on every tier. Easy through Med-C
retain large positive slack for both methods ($4.8$--$8.4$).
Hard-A: ME reaches $2.28$ while DE only achieves $0.58$, both
above zero but ME comfortably so. Hard-B: ME hits the floor at
$0.00$ while DE goes \emph{negative} ($-1.89$) --- DE cannot find a
layout with positive pairwise margin net of the variance penalty
under the Bonferroni-corrected bound for $M{=}8$ goals at the
expected horizon.

\paragraph{Argmax accuracy.}
Both DE and ME achieve $100\%$ argmax on Easy through Med-C,
matching theoretical predictions from their positive slacks.  On
Hard-A ($M{=}5$), ME reaches $\mathbf{100\%}$ (slack $2.28$) while
DE reaches only $94\%$ (slack $0.58$); the gap matches expectation
--- DE's marginal slack is insufficient to robustly guarantee
argmax under curved trajectories, while ME's larger slack clears
the bound comfortably. On Hard-B ($M{=}8$) the Boltzmann observer
reaches its discrimination limit: ME argmax drops to $54\%$ and DE
to $22\%$, with random at $25\%$. DE's argmax is actually
\emph{below} random here because DE's negative slack indicates no
disambiguating geometry exists under the bound --- DE's single-point
optimization converges on a locally optimal but globally uninformative
layout. ME's archive search still finds layouts that more than
double random's argmax on Hard-B ($54\%$ vs $25\%$), despite
zero-slack in the bound.

\paragraph{Time to inference.}
DE is the fastest on Easy ($4.1$\,s), Med-B ($8.1$\,s), and Med-C
($3.6$\,s); ME wins Hard-A ($10.6$ vs $18.0$\,s) and Hard-B
($9.6$ vs $17.6$\,s). On the Hard tiers DE runs close to the
$30$\,s timeout because its layouts produce indecisive posteriors
that never cross the $0.9$ commit threshold, relying on the $2$\,cm
arrival rule to exit. ME's Hard-tier layouts produce more
concentrated posteriors, hitting arrival sooner. Random is faster
than ME on Med-A ($6.9$ vs $9.4$\,s), partly because random's $33\%$
infeasibility excludes the hardest picks from its timed subset.

\paragraph{Grasp feasibility.}
A critical practical advantage of the optimized layouts is grasp
feasibility.  Random layouts have $5$--$49\%$ per-pick infeasibility
(arm-vs-other-object collision at the IK-solved grasp pose), which
translates to task success rates as low as $10\%$ on Hard-A. Both
DE and ME produce layouts with $0\%$ infeasibility and $100\%$
task success on every tier.  The optimizer's collision-aware feasibility
check---which walks MuJoCo contact pairs at each IK solution---ensures
that every grasp in the optimized layout can be physically executed
without the arm sweeping through neighboring objects.

\paragraph{Seed sensitivity.}
Slack is a noisy predictor of true argmax accuracy in the
approximation-limited regime.  On Hard-A, four independent ME runs
(seeds $42$--$45$) produced slacks in $[2.28, 3.20]$ with a
$30$-point spread in argmax accuracy ($70$--$100\%$), and the
\emph{highest-slack} seed ($3.20$) yielded only $80\%$ argmax
while the \emph{lowest-slack} seed ($2.28$) achieved $100\%$.
The ME results reported in \cref{tab:main_results} use the seed-$45$
layout, which exhibits the cleanest geometry under curved trajectories
despite having slightly lower nominal slack.  For deployment we
recommend either (i) selecting the best elite from the ME archive by
direct SA-in-the-loop evaluation, or (ii) tightening the slack
objective to penalize near-collinear goal pairs relative to the
start pose.

% ------------------------------------------------------------------
\subsection{RQ2: Scaling with the Number of Goals}
\label{sec:rq2_results}

\paragraph{Inference performance vs.\ $M$.}
\cref{fig:scaling} summarizes how argmax accuracy and time-to-inference
scale with $M$ across conditions.  For ME-optimized layouts, argmax
accuracy remains at $97$--$100\%$ from $M{=}2$ through $M{=}8$,
indicating that the optimizer consistently finds layouts where the
observer identifies the correct goal.  Time-to-inference grows from
$\sim 3$\,s at $M{=}2$ to $\sim 15$\,s at $M{=}8$, primarily due
to increased EE travel distance (the optimizer spreads objects wider
to maintain angular separation) and Boltzmann saturation delaying
threshold crossing.

For random baselines, argmax accuracy degrades from $88$\% to $67$\%
as $M$ increases, and task success rate collapses from $90$\% to
$23$\%---reflecting that both inference quality and physical
feasibility worsen as the workspace fills up.

\paragraph{Threshold sensitivity.}
We sweep $p_{\mathrm{thresh}} \in \{0.5, 0.6, 0.7, 0.8, 0.85, 0.9,
0.95\}$ using a bypass-threshold mode that records the first
crossing step for each $p_{\mathrm{thresh}}$ in a single run.
\cref{tab:sweep} reports the results for ME-optimized layouts.
Easy and Med-A achieve $100$\% threshold accuracy at every
$p_{\mathrm{thresh}}$ from $0.6$ to $0.9$, dropping only at
$p_{\mathrm{thresh}}{=}0.95$.  Med-B reaches $100$\% from $0.6$ to
$0.9$.  Med-C is stable at $94$\% from $0.8$ to $0.95$.
Hard-A peaks at $80$\% (${p_{\mathrm{thresh}}}{=}0.8$) then drops
to $44$\% at $0.9$, and Hard-B drops from $84$\% at $0.5$ to
$25$\% at $0.9$, consistent with the Boltzmann saturation bound:
for $M{=}8$ alternatives at $\beta{=}5$, the softmax posterior
physically cannot exceed $\sim 0.91$ regardless of layout.

Importantly, argmax accuracy remains $97$--$100\%$ on all tiers
regardless of threshold---the cost ordering is correct even when
the posterior is not concentrated enough to cross a fixed threshold.
This suggests that an argmax-with-margin decision rule would
outperform a fixed-threshold rule on high-$M$ scenes.

At low thresholds ($p_{\mathrm{thresh}} {\le} 0.5$), ME and random
baselines both suffer from premature wrong commits: the posterior at
step~1 is close to uniform, and whichever goal happens to have
marginally higher noisy posterior gets committed on.  This drops
threshold accuracy to $50$--$67$\% even for ME on easy scenes.
Above $p_{\mathrm{thresh}} {=} 0.7$, premature wrong commits
are effectively zero for ME on every tier.

\begin{table}[t]
\centering
\caption{Threshold sensitivity for ME-optimized layouts.  Each cell
is the threshold accuracy (fraction of picks where the posterior
first crossed $p_{\mathrm{thresh}}$ on the correct goal).  The sweep
is run with bypass-threshold mode.}
\label{tab:sweep}
\small
\setlength{\tabcolsep}{4pt}
\begin{tabular}{@{}lccccccc@{}}
\toprule
\textbf{Tier} &
  $0.5$ & $0.6$ & $0.7$ & $0.8$ & $0.85$ & $0.9$ & $0.95$ \\
\midrule
Easy   &  50\% &  50\% & 100\% & 100\% & 100\% & 100\% & 100\% \\
Med-A  & 100\% & 100\% & 100\% & 100\% & 100\% & 100\% &  67\% \\
Med-B  &  67\% & 100\% & 100\% & 100\% & 100\% & 100\% &  67\% \\
Med-C  & 100\% & 100\% & 100\% &  94\% &  94\% &  94\% &  94\% \\
Hard-A &  66\% &  60\% &  76\% &  80\% &  80\% &  44\% &  20\% \\
Hard-B &  84\% &  71\% &  62\% &  49\% &  37\% &  25\% &  12\% \\
\bottomrule
\end{tabular}
\end{table}

\paragraph{Optimizer compute cost.}
\cref{tab:compute} reports the optimization wall-clock time.  DE
runs in $30$--$92$\,s per scene, producing a single solution.  ME
runs in $340$--$920$\,s per scene ($\sim 10\times$ slower) but
produces a diverse archive of $374$--$400$ feasible layouts.  The
per-evaluation cost is identical for both (same yaw GD + IK +
collision pipeline); the difference is purely in the number of
candidates evaluated ($\sim 15$K for DE vs.\ $\sim 55$K--$90$K
for ME).  Both are offline costs acceptable for workspace design.

\begin{table}[t]
\centering
\caption{Optimizer compute cost (wall-clock seconds, single seed).
  ME additionally reports archive diversity metrics.}
\label{tab:compute}
\small
\begin{tabular}{@{}lccccc@{}}
\toprule
\textbf{Tier} & $M$ &
  DE (s) & ME (s) &
  Elites & Coverage \\
\midrule
Easy   & 2 &  55 &  597 & 374 & 93.5\% \\
Med-A  & 3 &  77 &  920 & 399 & 99.8\% \\
Med-B  & 3 &  34 &  653 & 400 & 100\%  \\
Med-C  & 3 &  92 &  727 & 397 & 99.3\% \\
Hard-A & 5 &  30 &  461 & 400 & 100\%  \\
Hard-B & 8 &  56 &  340 & 400 & 100\%  \\
\bottomrule
\end{tabular}
\end{table}

The archive diversity is a practical advantage of ME over DE.
With $93$--$100$\% coverage of the behavior space, the archive
provides hundreds of geometrically distinct feasible layouts.
This enabled the Med-A result: DE's single solution had poor
path-efficiency discrimination ($8$\% argmax), but ME's archive
contained a layout that achieves $100$\% on the same evaluation
observer.  A practitioner can evaluate the top-$K$ elites against
their deployment observer and select the best, hedging against
potential mismatch between the optimization objective and the
runtime inference model.


\end{document}
